\documentclass[10pt]{article}
\usepackage[T1]{fontenc}
\usepackage[utf8]{inputenc}
\usepackage[french]{babel}
\usepackage{amsmath,amssymb,mathrsfs,mathpazo}
\usepackage{geometry}
\geometry{a4paper,margin=2cm}
\usepackage{graphicx,array,tabularx,tabularray,colortbl,multirow,multicol,booktabs,enumitem,pifont,ulem,fancybox,diagbox,cancel}
\usepackage{pstricks,pstricks-add,pst-plot,pst-tree,pst-node,pst-text,pst-func,pst-math,pst-bspline,pst-3dplot,pst-eucl,pst-solides3d}
\let\pscalculate\undefined
\usepackage{pst-bezier}
\usepackage{tikz,pgfplots,tkz-tab}
\usetikzlibrary{arrows,arrows.meta,patterns,calc,positioning,trees,decorations.markings,shapes.geometric,shapes.misc,shapes.arrows,intersections,decorations.pathmorphing,decorations.pathreplacing,backgrounds,shadows,babel}
\pgfplotsset{compat=1.18}
\usepgfplotslibrary{fillbetween}
\usepackage[np]{numprint}
\usepackage{listings,makecell,pas-tableur,twoopt}
\let\barre\undefined
\usepackage[pstricks]{tablvar}
\definecolor{darkgreen}{RGB}{0,100,0}
\providecommand{\red}{\color{red}}
\providecommand{\blue}{\color{blue}}
\providecommand{\green}{\color{green}}
\providecommand{\black}{\color{black}}
\providecommand{\white}{\color{white}}
\providecommand{\gray}{\color{gray}}
\providecommand{\cyan}{\color{cyan}}
\providecommand{\magenta}{\color{magenta}}
\providecommand{\yellow}{\color{yellow}}
\providecommand{\violet}{\color{violet}}
\providecommand{\orange}{\color{orange}}
\providecommand{\euro}{\texteuro}
\providecommand{\EUR}{\texteuro}
\providecommand{\ieme}{\textsuperscript{e}}
\providecommand{\degres}{\ensuremath{^\circ}}
\providecommand{\degree}{\ensuremath{^\circ}}
\providecommand{\vv}[1]{\overrightarrow{#1}}
\providecommand{\Coord}[1]{#1}
\providecommand{\decosix}{}
\providecommand{\starredbullet}{\ensuremath{\bullet}}
\pagestyle{empty}
\setlength{\parindent}{0pt}
\usepackage{adjustbox}\DeclareUnicodeCharacter{2236}{:}\DeclareUnicodeCharacter{2190}{\ensuremath{\leftarrow}}
\providecommand{\R}{\mathbb{R}}
\providecommand{\N}{\mathbb{N}}
\providecommand{\Z}{\mathbb{Z}}
\providecommand{\Q}{\mathbb{Q}}
\providecommand{\D}{\mathbb{D}}
\providecommand{\C}{\mathbb{C}}
\providecommand{\e}{{\mathrm{e}}}
\providecommand{\vect}[1]{\overrightarrow{#1}}
\providecommand{\vectt}[1]{\overrightarrow{\mathrm{#1}}}
\providecommand{\vvt}[1]{\overrightarrow{\mathrm{#1}}}
\providecommand{\np}[2][]{{#2~#1}}
\providecommand{\vv}[1]{{\overrightarrow{#1}}}
\renewcommand{\d}{{\mathrm{d}}}
\providecommand{\degres}{{^\circ}}
\providecommand{\degre}{{^\circ}}
\providecommand{\blue}{\color{blue}}
\providecommand{\red}{\color{red}}
\providecommand{\enskip}{\quad}
\providecommand{\llbracket}{\lbrack\!\lbrack}
\providecommand{\rrbracket}{\rbrack\!\rbrack}
\providecommand{\euro}{{\text{€}}}
\providecommand{\no}{{\text{n°}}}
\providecommand{\interval}[2]{{[#1;#2]}}
\providecommand{\intervFO}[2]{{[#1;#2[}}
\providecommand{\dotfill}{\ldots}

\providecommand{\ds}{\displaystyle}
\providecommand{\textding}[1]{#1}
\providecommand{\Oij}{\(\mathrm{O};\vec{\imath},\vec{\jmath}\)}
\providecommand{\Oijk}{\(\mathrm{O};\vec{\imath},\vec{\jmath},\vec{k}\)}
\providecommand{\Ouv}{\(\mathrm{O};\vec{u},\vec{v}\)}
\providecommand{\intervFF}[2]{\ensuremath{[#1;#2]}}
\renewcommand{\euro}{\ifmmode\text{\texteuro}\else\texteuro\fi}
\DeclareUnicodeCharacter{00B0}{\ifmmode^{\circ}\else\textdegree\fi}
\DeclareUnicodeCharacter{20AC}{\ifmmode\text{\texteuro}\else\texteuro\fi}
% Énoncés : Éducation nationale. Transcription et corrigés : APMEP, Denis Vergès et contributeurs. Découpage et classement : CoopMaths, Sébastien Lozano, Éric Elter et contributeurs. Adaptation de la présentation : Maths Explorer.
% CC BY-SA : https://coopmaths.fr/www/about/
% https://coopmaths.fr/alea/static/bac/2025/tex/bac_2025_09_sujet2_metropole_3_cor.tex
\begin{document}
\noindent\textit{APMEP / CoopMaths. CC BY-SA. Présentation adaptée par Maths Explorer.}\par\medskip
\subsection*{Partie A}

\begin{enumerate}
	\item $-x^2 + 7x + 8 =0$ admet $-1$ comme solution évidente. L'autre solution de l'équation du second degré est donc $8$.

Le trinôme est négatif (du signe du coefficient de $x^2$), sauf entre les racines, d'où le tableau de signes  de $x\longmapsto -x^2 + 7x + 8 $ :
	
	\begin{tikzpicture}
	\tkzTabInit[lgt=5.5,espcl=2]
	{$x$
	/1,
	signe de $-x^2+7x+8$
	/1
	}
	{$-\infty$ , $-1$ , $8$ , $+\infty$}
	\tkzTabLine{, - , z , + , z, -, }
	\end{tikzpicture}
	
L'ensemble des solutions de l'inéquation 	$-x^2 + 7x + 8 \geqslant 0$ sur $\R$ est l'intervalle $\left[-1~;~8\right]$
	
	\item Pour tout nombre réel $x$ appartenant à l'intervalle $\left]0~;~8\right]$, $-x^2 + 7x + 8 \geqslant 0$. D'où en ajoutant 1 à chaque membre $-x^2 + 7x + 9 \geqslant 1$.

	Comme la fonction logarithme népérien est une fonction croissante on a pour tout  $x$ appartenant à $\left]0~;~8\right]$, $\ln\left(- x^2 + 7x + 9\right) \geqslant \ln(1)$. C'est à dire $\ln\left(- x^2 + 7x + 9\right) \geqslant 0$.

Pour tout $x$ appartenant à $\left]0~;~8\right]$, $x>0$ et  $10\ln\left(-x^2 + 7x + 9\right) \geqslant 0$, donc $f(x) = \dfrac{10\ln\left(-x^2 + 7x + 9\right)}{x} \geqslant 0$.

	\item Sur l'intervalle $[0~;~8]$ la courbe représentative est au dessus de l'axe des abscisses.
	
\end{enumerate}

\subsection*{Partie B}

\begin{enumerate}
	\item Pour $x\in \left]0~;~8\right]$, M$\left(x,\dfrac{10\ln\left(-x^2 + 7x + 9\right)}{x}\right)$ et donc N$\left(x,0\right)$ et P$\left(0,\dfrac{10\ln\left(-x^2 + 7x + 9\right)}{x}\right)$.
	\item Comme pour $x\in \left]0~;~8\right]$, l'abscisse de N est positive et l'ordonnée de P est positive. On peut écrire ON$=x$ et OP$=\dfrac{10\ln\left(-x^2 + 7x + 9\right)}{x}$.
	
	L'aire de ONMP vaut $\text{ON}\times\text{OP}$ donc : 
	
	$\forall x \in \left]0~;~8\right]\quad \mathcal{A}(x) = x\times \dfrac{10\ln\left(-x^2 + 7x + 9\right)}{x} = 10\ln\left(-x^2 + 7x + 9\right)$
	\item Sur $\left]0~;~8\right[$, $-x^2 + 7x + 9 > 0$ donc $\mathcal{A}$ est une fonction dérivable de la variable $x$ sur $\left]0~;~8\right[$ comme composée de fonction définies et dérivables. 
	
	$\forall x\in \left]0~;~8\right[$, $\mathcal{A}'(x) = 10\times \dfrac{-2x + 7}{-x^2 + 7x +9}$.
	
	S'il existe $x\in \left]0~;~8\right[$ telle que $\mathcal{A}(x)$ soit maximale, il faut que $\mathcal{A}'(x) = 0$
	
	Pour $x\in \left]0~;~8\right[$, $-x^2 + 7x +9 > 0 $,  $\mathcal{A}'(x)$ est donc du signe de $-2x + 7$.
	
	$
	\begin{array}{ccc}
	\mathcal{A}'(x) \geqslant 0 &\iff& -2x + 7 \geqslant 0 \\
	\mathcal{A}'(x) \geqslant 0 &\iff& x \leqslant \dfrac{7}{2} \\
	\end{array}
	$

	$\mathcal{A}$ est croissante sur $\left]0~;~\dfrac{7}{2}\right[$ et décroissante sur  $\left]\dfrac{7}{2}~;~8\right[$.
	
	La dérivée s'annule en $\dfrac{7}{2}$, en étant positive avant puis négative après, donc la fonction $\mathcal{A}$ a un maximum sur $]0~;~8[$, : 
	
	$\mathcal{A}\left(\dfrac{7}{2}\right) = 10\ln(17)+ 10\ln(5)-20\ln(2) \approx \np{30.56}$
	
	\end{enumerate}
	
\subsection*{Partie C}

\begin{enumerate}
	\item \ {} 
	
\begin{center}
\begin{lstlisting}[language=python,frame=single, escapechar=|, numbers=left]
from math import *	

def A(x):
	return 10*log(-1 * x**2 + 7*x + 9)
	
def pluspetitevaleur(k):
	x = 3.5
	while A(x) > k:
		x = x + 0.1
	return x	

\end{lstlisting}
\end{center}
	
	\item
	\begin{lstlisting}[language=python]
>>>pluspetitevaleur(30)
4.599999999999998
\end{lstlisting}
	
	\item Lorsque $k=35$, la condition \verb!A(x) > 35! n'est pas vérifiée donc le contenu de la boucle n'est pas exécuté. L'appel \verb!pluspetitevaleur(35)! renvoie \verb!3.5!.
\end{enumerate}

\end{document}
